\documentclass[11pt]{article}
\usepackage[margin=1in]{geometry}
\usepackage{amsmath}

\setlength{\parindent}{0pt}

\begin{document}

\textbf{Problem 3.1}
\bigskip

The assets of Dallas \& Associates consist entirely of current assets and net plant and equipment, and the firm has no excess cash. The firm has total assets of \$2.5 million and net plant and equipment equals \$2 million. It has notes payable of \$150,000, long-term debt of \$750,000, and total common equity of \$1.5 million. The firm does have accounts payable and accruals on its balance sheet. The firm only finances with debt and common equity, so it has no preferred stock on its balance sheet.

\bigskip
\textbf{Answer}

\begin{enumerate}
  \item[a.] \textbf{Total debt:}
    \[
      \text{Total Debt} = \text{Notes Payable} + \text{Long-Term Debt} = 150{,}000 + 750{,}000 = \boxed{\$900{,}000}
    \]

  \item[b.] \textbf{Total liabilities and equity:}

    By the accounting identity, total liabilities and equity equals total assets:
    \[
      \boxed{\$2{,}500{,}000}
    \]

  \item[c.] \textbf{Current assets:}
    \[
      \text{Current Assets} = \text{Total Assets} - \text{Net P\&E} = 2{,}500{,}000 - 2{,}000{,}000 = \boxed{\$500{,}000}
    \]

  \item[d.] \textbf{Current liabilities:}

    First, find accounts payable and accruals (from part e below):
    \[
      \text{Total Liabilities} = \text{Total Assets} - \text{Total Equity} = 2{,}500{,}000 - 1{,}500{,}000 = 1{,}000{,}000
    \]
    \begin{align*}
      \text{AP \& Accruals} &= \text{Total Liabilities} - \text{Notes Payable} - \text{LTD} \\
                            &= 1{,}000{,}000 - 150{,}000 - 750{,}000 = 100{,}000
    \end{align*}
    \[
      \text{Current Liabilities} = \text{Notes Payable} + \text{AP \& Accruals} = 150{,}000 + 100{,}000 = \boxed{\$250{,}000}
    \]

  \item[e.] \textbf{Accounts payable and accruals:}
    \begin{align*}
      \text{AP \& Accruals} &= \text{Total Liabilities} - \text{Notes Payable} - \text{LTD} \\
                            &= 1{,}000{,}000 - 150{,}000 - 750{,}000 = \boxed{\$100{,}000}
    \end{align*}

  \item[f.] \textbf{Net working capital (NWC):}
    \[
      \text{NWC} = \text{Current Assets} - \text{Current Liabilities} = 500{,}000 - 250{,}000 = \boxed{\$250{,}000}
    \]

  \item[g.] \textbf{Net operating working capital (NOWC):}

    Operating current assets exclude excess cash (none here). Operating current liabilities exclude interest-bearing debt (notes payable):
    \[
      \text{NOWC} = \text{Operating CA} - \text{Operating CL} = 500{,}000 - 100{,}000 = \boxed{\$400{,}000}
    \]

  \item[h.] \textbf{Explanation of difference:}

    NWC (\$250,000) is less than NOWC (\$400,000) by exactly \$150,000, which equals the notes payable. NWC subtracts all current liabilities including notes payable (an interest-bearing, financing liability). NOWC excludes notes payable because it is a financing source, not an operating obligation. Only accounts payable and accruals---which arise from operations---are treated as operating current liabilities in NOWC.
\end{enumerate}

\bigskip
\bigskip

\textbf{Question 3.9}
\bigskip

Which of the following actions are most likely to directly increase cash as shown on a firm's balance sheet?

\bigskip
\textbf{Answer}

\begin{enumerate}
  \item[a.] \textbf{It issues \$4 million of new common stock.} \textit{Yes, this directly increases cash.} When a firm issues equity, it receives cash proceeds from investors. Cash rises by \$4 million and stockholders' equity rises by the same amount.

  \item[b.] \textbf{It buys new plant and equipment at a cost of \$3 million.} \textit{No, this decreases cash} (assuming cash purchase) or increases debt. Acquiring a long-term asset is a use of funds, not a source.

  \item[c.] \textbf{It reports a large loss for the year.} \textit{No, not directly.} A reported loss is an accounting entry reflecting revenues minus expenses. It does not itself generate or consume cash; the cash impact depends on the nature of the underlying transactions (e.g., non-cash charges like depreciation do not reduce cash).

  \item[d.] \textbf{It increases the dividends paid on its common stock.} \textit{No, this decreases cash.} Paying dividends is a cash outflow to shareholders and reduces the cash balance.
\end{enumerate}

\end{document}
